\clearpage\fancyhead[L]{\small JUNO.}\hypertarget{native-365-lower}{}\label{native-365-lower}
\begin{center}\textbf{JUNO.}\par\vspace{4pt}\hrule\vspace{4pt}\fontsize{10}{12}\selectfont Monatliche Correspondenz. Vol. XXVI. pp. 297\;..\;299. September 1812.\par\vspace{4pt}\hrule\vspace{4pt}\end{center}

\begin{flushright}Göttingen, 9 Sept. 1812.\end{flushright}
\par\indent Today I want to follow up with my observations of \emph{Juno}. Finding this small planet took an extraordinary amount of effort this year, since it barely reached the 11th magnitude and lay right in the Milky Way. On top of that, moonlight and poor weather interrupted the search. I therefore obtained only four observations, though I consider all of them very good.

\begin{center}\fontsize{10}{12}\selectfont\setlength{\tabcolsep}{4pt}\renewcommand{\arraystretch}{1.15}\begin{tabular}{l|c|c|c}
1812 & Mean time & R.A. & Declination\\\hline
June 7 & $12^{\mathrm u}\;21^{\mathrm m}\;0^{\mathrm s}$ & $282^\circ\;21'\;24''5$ & $4^\circ\;41'\;44''6$ S.\\
13 & $12\;51\;12$ & $281\;16\;32.9$ & $4\;33\;39.6$\\
July 12 & $11\;35\;28$ & $275\;5\;10.3$ & $5\;6\;11.2$\\
13 & $11\;28\;12$ & $274\;52\;49.9$ & $5\;9\;20.0$\\
\end{tabular}\end{center}

\Needspace{13\baselineskip}\par\indent My X elements give the following differences:

\begin{center}\fontsize{10}{12}\selectfont\setlength{\tabcolsep}{7pt}\renewcommand{\arraystretch}{1.15}\begin{tabular}{l|c|c}
& R.A. & Decl.\\\hline
June 7 & $+21''8$ & $-35''9$\\
13 & $+31.0$ & $-34.7$\\
July 12 & $+22\;4$ & $-46.9$\\
13 & $+26.2$ & $-45.0$\\
\end{tabular}\end{center}

\par\indent Mr. \textsc{Wachter} took on the task of calculating the opposition from these observations. I recommended a method I have long used for this purpose. It has the advantage that observations somewhat farther away in time can also be used for the opposition without concern. I will write to you another time about this procedure in a little more detail. Mr. \textsc{Wachter}’s result is as follows:

\Needspace{8\baselineskip}\begin{center}
\emph{Opposition of Juno} 1812.
\end{center}

29 June $19^{\mathrm u}\;51^{\mathrm m}\;43^{\mathrm s}$ Mean time in Göttingen

\begin{center}\fontsize{10}{12}\selectfont\setlength{\tabcolsep}{2pt}\renewcommand{\arraystretch}{1.15}\begin{tabular}{lcr}
True longitude & $.\;.\;.\;.\;.\;.\;.\;.\;.\;.\;.$ & $278^\circ\;15'\;21''3$\\
True geocentric latitude & $.\;.\;.\;.\;.\;.$ & $18\;36\;56.9$\\
\end{tabular}\end{center}

\Needspace{19\baselineskip}\par\indent Mr. \textsc{Wachter} also corrected the elements on the basis of the last four observed oppositions (the one in 1810 was missed, as you know). He followed the method presented in volume I of the \emph{Comment. Nov. Soc. Sc. Götting.} [p. 1 of this volume] and found the following:
\SourceLeaf{366}\begin{center}\fontsize{9}{11}\selectfont\setlength{\tabcolsep}{2pt}\renewcommand{\arraystretch}{1.15}\begin{tabular}{lcr}
Epoch 1813, Göttingen meridian & $.\;.\;.\;.\;.\;.\;.$ & $342^\circ\;56'\;32''2$\\
Daily mean tropical motion & $.\;.\;.\;.\;.\;.$ & $813''25748$\\
Perihelion 1813 & $.\;.\;.\;.\;.\;.\;.\;.\;.\;.\;.\;.\;.\;.$ & $53^\circ\;16'\;52''3$\\
Node 1813 & $.\;.\;.\;.\;.\;.\;.\;.\;.\;.\;.\;.\;.\;.\;.\;.$ & $171\;10\;58.9$\\
Inclination of the orbit & $.\;.\;.\;.\;.\;.\;.\;.\;.\;.\;.\;.\;.$ & $13\;4\;17.2$\\
Angle of eccentricity & $.\;.\;.\;.\;.\;.\;.\;.\;.\;.\;.\;.$ & $14\;43\;58.1$\\
Logarithm of the semimajor axis & $.\;.\;.\;.\;.\;.\;.$ & $0.4265679$\\
\end{tabular}\end{center}

\par\indent Now that the secular changes in the elements of the orbits of Ceres and Pallas are known insofar as they are due to Jupiter (compared with whose effect the other planets are quite insignificant), I asked Mr. \textsc{Enke}, a talented young man and a calculator as skilled as he is careful, to investigate whether the distance between the orbits of Ceres and Pallas, which is known to be fairly small at the ascending node of the orbit of Ceres on that of Pallas — the very circumstance that led Dr. \textsc{Olbers} to his well-known hypothesis — is now decreasing or increasing, that is, whether the orbits are approaching an actual intersection or moving away from one. The latter would obviously favor \textsc{Olbers}’s hypothesis, but the result is just the reverse. The distance was formerly greater than it is now. Mr. \textsc{Enke} finds the following radius vectors

\begin{center}\fontsize{9}{11}\selectfont\setlength{\tabcolsep}{2pt}\renewcommand{\arraystretch}{1.15}\begin{tabular}{c|c|c|c||c|c|c}
& \multicolumn{3}{c||}{\shortstack{at the $\AscendingNode$ of the orbit of Ceres\\on that of Pallas}} & \multicolumn{3}{c}{\shortstack{at the $\DescendingNode$ of the orbit of Ceres\\on that of Pallas}}\\
Year & Ceres & Pallas & Difference & Ceres & Pallas & Difference\\\hline
$808$ & $2.82294$ & $2.70322$ & $-0.11972$ & $2.67612$ & $2.37780$ & $-0.29832$\\
$1808$ & $2.92427$ & $2.84945$ & $-0.07482$ & $2.59569$ & $2.40346$ & $-0.19223$\\
$3475$ & $2.95374$ & $2.95743$ & $+0.00369$ & $2.57987$ & $2.49163$ & $-0.08824$\\
\end{tabular}\end{center}

\par\indent On this basis, an actual intersection at the $\AscendingNode$ would occur around the year 3397, which may be regarded as approximately correct. Of course, there will also have been an intersection at some time in the past. From the trend of the numbers in the fourth and sixth columns, however, we can at least conclude that this could only have happened at an epoch many thousands of years away. Thus if one wishes to accept Dr. \textsc{Olbers}’s hypothesis about the origin of the new planets, that origin lies in a past so distant that it cannot yet be calculated, before the times covered by our history.
